% !TeX root = ../测度论笔记.tex
% 由主文档通过 \input 载入；请编译主文档。
\setcounter{chapter}{11} % 保留原书章号。
\chapter{积分与微分的联系：变差、密度与绝对连续性}
\label{chap:mt-calculus}

\section{为什么需要重新审视微积分基本定理}
数学分析中，连续函数的变上限积分可微，导数等于被积函数；连续可微函数又能由其导数积分恢复。这两个方向在连续性条件下十分整齐。进入测度论后，被积函数只要求可积，原函数可能只在几乎处处意义下可微，因而必须重新说明结论中的量词与假设。

首先应分开三个问题：一个函数是否几乎处处存在导数，导数是否可积，以及导数的积分能否恢复原函数。第一个问题有肯定答案，并不自动给出后两个问题的答案。即使导数几乎处处存在且等于零，函数也未必为常数；额外需要排除集中在零测集上的变化。绝对连续性正是完成这一排除的适当条件。

本章作为测度论与后续分析课程之间的补充路线。基础学习可先掌握定义、典型反例与积分表示，再在进一步的实分析课程中研究微分定理的覆盖论证。它不影响前面已完成的抽象积分理论，但能说明勒贝格积分如何重新连接熟悉的微积分。

\section{有界变差函数与 Stieltjes 测度}
\begin{definition}
对 $F:[a,b]\to\mathbb R$，定义总变差
\[
 V_a^b(F)=\sup_{a=x_0<\cdots<x_m=b}
 \sum_{j=1}^{m}|F(x_j)-F(x_{j-1})|.
\]
若该上确界有限，称 $F$ 为\emph{有界变差函数（function of bounded variation）}。
\end{definition}
这里的变差累计所有上升与下降，而端点差只记录净变化。单调函数的变差等于端点差的绝对值，连续可微函数的变差不超过 $\int_a^b|F'|$；频繁振荡可能使变差很大，即使函数值始终处于一个很小的范围内。

若 $F$ 在实轴上非减且右连续，则可以在半开区间上规定
$\mu_F((u,v])=F(v)-F(u)$。半开区间可以无歧义地拼接，相邻区间的端点不会被重复计数；右连续性保证从上缩小区间时的相容性。由相应预测度扩张，可得到局部有限的 Borel 测度，称为\emph{Lebesgue--Stieltjes 测度（Lebesgue--Stieltjes measure）}。这里引用非减右连续函数所定义区间预测度的标准构造，完整验证见\cite[第1章]{teschl-ra}。

对单点，由区间 $(x-h,x]$ 从上趋于 $\{x\}$，可得
$\mu_F(\{x\})=F(x)-F(x-)$，其中 $F(x-)$ 是左极限。因此函数的跳跃对应测度的原子。若 $F$ 连续，单点质量都为零，但测度未必相对于长度绝对连续；连续性只排除了跳跃，尚未排除集中在不可数零测集上的质量。

\section{绝对连续性比一致连续性更强}
\begin{definition}
函数 $F:[a,b]\to\mathbb R$ 称为\emph{绝对连续的（absolutely continuous）}，若对每个 $\varepsilon>0$，存在 $\delta>0$，使任意有限个两两不交的子区间 $(a_j,b_j)\subseteq[a,b]$ 只要满足 $\sum_j(b_j-a_j)<\delta$，就有
$\sum_j|F(b_j)-F(a_j)|<\varepsilon$。
\end{definition}
一致连续性控制一段很短区间上的函数增量；绝对连续性同时控制任意多段小区间的增量总和。后者当然蕴含前者，只需取一个区间；反向则一般不成立，因为很多小变化可能在数量增加时累积成不可忽略的总量。

每个 Lipschitz 函数都绝对连续。若 $|F(x)-F(y)|\leq L|x-y|$，对不交区间求和便有总增量不超过 $L$ 乘总长度，可以取 $\delta=\varepsilon/L$，常函数情形另行直接处理。但绝对连续函数未必 Lipschitz，例如 $F(x)=\sqrt{x}$ 在 $[0,1]$ 上绝对连续，其导数 $1/(2\sqrt{x})$ 可积，却在零点附近无界。

绝对连续函数也有界变差。取对应于误差一的长度阈值 $\delta$，将 $[a,b]$ 预先分成有限个长度小于 $\delta$ 的子区间。对任意分割加入这些固定端点，变差和只会增大；每个固定子区间内所有小段的总长度小于 $\delta$，故其增量绝对值之和小于一。对有限个固定子区间求和，就得到与原分割无关的变差上界。

\section{可积函数的变上限积分}
\begin{proposition}
若 $f\in L^1([a,b])$，令 $F(x)=\int_a^xf(t)\,\mathrm dt$，则 $F$ 绝对连续。
\label{prop:mt-integral-ac}
\end{proposition}
\begin{proof}
对有限个两两不交区间，
\[
 \sum_j|F(b_j)-F(a_j)|
 \leq\sum_j\int_{a_j}^{b_j}|f(t)|\,\mathrm dt
 =\int_{\bigcup_j(a_j,b_j)}|f(t)|\,\mathrm dt.
\]
可积函数的积分对集合测度绝对连续，而该并集的长度就是各区间长度之和，因此满足函数绝对连续性的定义。
\end{proof}
这一结论并不要求 $f$ 连续，也不要求处处有界。它首先保证原函数的整体变化受集合长度控制，随后才能进一步研究导数。特别地，若把 $f$ 在零测集上修改，所有区间积分都不变，所以由积分定义的 $F$ 完全不变，而不仅仅是几乎处处不变。

\begin{theorem}[勒贝格微分定理的一维形式，引用]
若 $f$ 在实轴上局部可积，即在每个有界区间上绝对可积，则对几乎每个 $x$，
\[
 \lim_{r\downarrow0}\frac1{2r}\int_{x-r}^{x+r}|f(t)-f(x)|\,\mathrm dt=0.
\]
这些点称为 $f$ 的勒贝格点。
\label{thm:mt-lebesgue-differentiation}
\end{theorem}
完整证明使用覆盖选择与极大函数估计，本章仅引用其准确形式，参见\cite[第22--23讲]{mit-measure}。从对称区间平均误差趋零，可以推出左右单侧平均误差也趋零，因为单侧误差不超过两倍对称平均误差。因此，对命题\ref{prop:mt-integral-ac}中的 $F$，差商
$[F(x+h)-F(x)]/h$ 在每个内部勒贝格点趋于 $f(x)$，故 $F'=f$ 几乎处处。

\begin{theorem}[绝对连续函数的积分表示，引用]
$F:[a,b]\to\mathbb R$ 绝对连续，当且仅当存在 $f\in L^1([a,b])$，使
$F(x)=F(a)+\int_a^xf(t)\,\mathrm dt$ 对每个 $x\in[a,b]$ 成立。此时 $F'=f$ 几乎处处，因而
$F(x)=F(a)+\int_a^xF'(t)\,\mathrm dt$。
\end{theorem}
从积分表示到绝对连续性及导数结论已经在上面说明；反向需要将绝对连续函数的增量构造成关于长度绝对连续的有符号测度，再用 Radon--Nikodym 定理得到密度，或使用等价的实分析证明。这里给出证明路线而不省略为“由求导即得”，因为一般绝对连续函数并没有处处连续的导数。

\section{前往偏微分方程：弱导数的初步概念}
经典导数逐点比较函数值，因而会对零测集上的修改十分敏感；但 $L^p$ 元素本身只在几乎处处意义下确定。若希望在可积函数空间中讨论微分，应寻找一种不依赖代表函数个别点值的定义。分部积分提供了办法：把导数转移到已知的光滑测试函数上，从而只让待研究函数出现在积分中。

\begin{definition}
设 $I$ 为开区间，$f,g$ 在 $I$ 上局部可积。若对每个光滑且紧支撑于 $I$ 的测试函数 $\varphi\in C_c^\infty(I)$，
\[
 \int_I f(x)\varphi'(x)\,\mathrm dx
 =-\int_I g(x)\varphi(x)\,\mathrm dx,
\]
则称 $g$ 为 $f$ 的\emph{弱导数（weak derivative）}。该定义中的等式对所有测试函数成立，并在积分意义下忽略零测修改。
\end{definition}
若 $f$ 连续可微，经典分部积分说明其通常导数就是弱导数，因为测试函数在区间边界附近为零，不产生边界项。弱导数允许更多函数进入微分理论。例如 $f(x)=|x|$ 在零点不可微，但分别在零点两侧分部积分，左右边界项抵消，得到弱导数为符号函数，零点处的取值任意。

并非每个局部可积函数都有局部可积的弱导数。对阶跃函数 $H=\mathbf1_{(0,\infty)}$，有 $\int H\varphi'=-\varphi(0)$。若存在局部可积 $g$ 使其成为弱导数，则应有 $\int g\varphi=\varphi(0)$。选择满足 $0\leq\varphi_n\leq1$、$\varphi_n(0)=1$、支撑缩到零点的光滑测试函数，局部可积性与控制收敛使左侧趋零，而右侧恒为一，矛盾。

该阶跃的微分作用仍可用 Dirac 测度描述，因为 $\int\varphi\,\mathrm d\delta_0=\varphi(0)$；只是这个导数不能由一个局部可积函数表示。这说明“函数导数”“测度形式的导数”和更一般的分布导数具有不同的对象范围。测度论在这里帮助区分了能够由密度表示的变化与集中在单点上的变化。

弱导数与磨光相配合，可以构造 Sobolev 空间和弱解理论。本书只建立入口：可积函数空间提供积分框架，测试函数将微分转成线性恒等式，范数估计保证近似过程有可控极限。完整的偏微分方程分析还需要同时控制函数及其弱导数，不能只依靠本章的单变量例子。

\section{Cantor 函数与奇异连续质量}
在 Cantor 集构造中，可以同时构造一个连续非减函数 $C:[0,1]\to[0,1]$：在每个被删除的中间三分之一区间上取常数，在保留区间上按二进制权重继续分配增长，并使 $C(0)=0$、$C(1)=1$。一种严格做法是逐步构造连续分段线性函数，每步在新删除区间上拉平，后续每个保留区间内的变化幅度不超过 $2^{-n}$。这些函数一致收敛，极限便连续且非减。

Cantor 集外的每个点都落在某个删除开区间内，而 $C$ 在该区间恒定，所以在那里导数为零。Cantor 集长度为零，因此 $C'=0$ 几乎处处；但 $C(1)-C(0)=1$，故导数积分无法恢复端点变化。这说明“连续且几乎处处可微”不足以替代绝对连续性，也说明不能把几乎处处成立的微分恒等式无条件积分为全局恒等式。

与 $C$ 对应的 Stieltjes 测度没有单点原子，因为 $C$ 连续；它在每个删除区间上的测度为零，因此全部质量集中于 Cantor 集。它既非离散测度，又没有相对于长度的密度，称为奇异连续测度。这样，原子、绝对连续与奇异连续三种质量形式都能在同一测度框架中得到解释。

\section{习题}
\begin{enumerate}
 \item 证明 Lipschitz 条件蕴含绝对连续，并用 $\sqrt{x}$ 说明反向不成立。
 \item 对阶跃函数 $F(x)=\mathbf1_{[0,\infty)}(x)$，计算其 Stieltjes 测度在各半开区间上的值，并识别对应测度。
 \item 若 $F(x)=\int_a^xf(t)\,\mathrm dt$ 且 $f=0$ 几乎处处，证明 $F$ 在每个点都等于零；再说明 Cantor 函数为什么不是这一命题的反例。
\end{enumerate}

